\scanpage{103}
a) b) c) avec l'axiome $\alpha$ équivaut à : la donnée de l'ens $\Gamma$ \inserted{$\ne\varnothing$}, et de la famille $(E_\alpha)_{\alpha\in t_0}$ de ses quotients
\[
\varphi_\alpha:\Gamma\longrightarrow E_\alpha,\qquad
\varphi:\Gamma\unsure{\text{---}|\longrightarrow}\prod_{\alpha\in t_0}E_\alpha,
\]
\deleted{et}\inserted{puis} des rev. d'ordre 2
\[
\widetilde E_\alpha\xrightarrow{p_\alpha} E_\alpha\qquad(\alpha\in t_0)
\]
et d'un \uline{syst. transitif} de\deleted{\unsure{d}} \inserted{$\Gamma$-}iso. entre les revêtements $\varphi_\alpha^*(\widetilde E_\alpha/E_\alpha)$ de $\Gamma$ (permettant de définir $\widetilde\Gamma$ par identification\ldots) Ceci posé, l'axiome $\beta_{\widetilde\Gamma})$ \deleted{implique}\inserted{équivaut à} l'axiome

$\beta_\Gamma$) Si $\{\alpha,\beta\}\in\unsure{\mathcal P|\mathfrak P}_2(t_0)$, alors
\[
\Gamma\xrightarrow{\varphi_{\alpha\beta}}E_\alpha\times E_\beta
\]
est \uline{bijectif}

La classification de ces ``complexes cubiques généralisés'' \deleted{\unsure{admet}} se décompose donc en deux questions exactement

I) Description des systèmes
\[
\left(\Gamma\xrightarrow{\varphi_\alpha}\Gamma_\alpha\right)_{\alpha\in t_0}
\]
de quotients (\unsure{paramétré|paramétrés} par $t_0$ de card 3) d'un ens $\Gamma\ne\varnothing$, satisfaisant l'axiome $\beta_\Gamma$
% archive sheet number: 102
